\chapter{Evaluación}
\label{cap:capitulo5}

\section{Sistema de evaluación de las prácticas}

Cada práctica se evalúa mediante una \textbf{rúbrica analítica} que combina criterios técnicos con un \textbf{Criterio Green IT} específico, alineado con el principio GSE trabajado en esa práctica. Las rúbricas tienen una doble función: \textit{sumativa} —permiten calificar el entregable— y \textit{formativa} —el alumnado las recibe al inicio de cada práctica y las usa como guía de autoevaluación antes de entregar—. Este enfoque conecta con la evaluación auténtica propia del ABP: no se evalúa solo el resultado final, sino el proceso de toma de decisiones técnicas y su justificación ambiental.

Los niveles de desempeño se calibran de forma que un alumno que resuelva correctamente el problema técnico sin prestar atención al impacto ambiental no supere el nivel \textit{Bien}; alcanzar \textit{Excelente} requiere siempre cuantificar, reflexionar o transferir el conocimiento green. Las rúbricas de P0–P4 se presentan a continuación.

\newpage
\FloatBarrier
\subsection{Rúbrica P0: \textit{¿Cuánto contamina tu código?}}

\begin{table}[htbp]
\centering
\scalebox{0.92}{\footnotesize
\begin{tabular}{|p{2.6cm}|p{2.6cm}|p{2.6cm}|p{2.6cm}|p{2.6cm}|}
\hline
\textbf{Criterio} & \textbf{Excelente (4)} & \textbf{Bien (3)} & \textbf{En desarrollo (2)} & \textbf{Insuficiente (1)} \\
\hline
\textbf{Comprensión del impacto TIC} & Describe con precisión al menos tres dimensiones del impacto del sector TIC apoyándose en datos cuantitativos contrastados & Describe dos dimensiones con datos cuantitativos correctos & Describe una dimensión sin datos cuantitativos o con datos imprecisos & No articula ninguna dimensión con datos concretos \\
\hline
\textbf{Participación en el Jigsaw} & Presenta su dominio con claridad, responde preguntas y completa la Ficha B íntegramente & Presenta su dominio con claridad pero la Ficha B queda incompleta & Presenta el dominio con apoyo docente o con datos escasos & No presenta su dominio o los datos son incorrectos \\
\hline
\textbf{Cálculo de huella personal} & Calcula correctamente la huella anual de datos móviles; traduce el resultado a equivalencias de la Tabla 4.2; cálculo documentado en la Ficha C & Cálculo correcto sin traducir a equivalencias & Cálculo parcialmente correcto (error en un factor) & No realiza el cálculo o los datos de partida no son reales \\
\hline
\textbf{Criterio Green IT} & Ficha C completa con huella cuantificada y reflexión sobre el resultado & Ficha C con cálculo correcto pero sin reflexión & Ficha C con cálculo incompleto & Ficha C en blanco o no entregada \\
\hline
\textbf{Reflexión crítica} & Identifica el sesgo de subestimación propio y explica al menos una causa estructural & Identifica el sesgo sin explicar la causa & Reconoce haber subestimado sin mayor elaboración & No participa en la reflexión \\
\hline
\end{tabular}}
\caption{Rúbrica P0: \textit{¿Cuánto contamina tu código?}}
\label{tab:rubrica-p0-cuanto-contamina-tu-codigo}
\end{table}

\newpage
\FloatBarrier
\subsection{Rúbrica P1: \textit{¿Cuánto pesa tu algoritmo?}}

\begin{table}[htbp]
\centering
\scalebox{0.92}{\scriptsize
\begin{tabular}{|p{2.4cm}|p{2.4cm}|p{2.4cm}|p{2.4cm}|p{2.4cm}|}
\hline
\textbf{Criterio} & \textbf{Excelente (4)} & \textbf{Bien (3)} & \textbf{En desarrollo (2)} & \textbf{Insuficiente (1)} \\
\hline
\textbf{Comprensión algorítmica} & Relaciona correctamente la complejidad $O(2^n)$ con el consumo energético antes de medir y predice cuál será más eficiente con justificación teórica & Identifica la diferencia de complejidad pero no la vincula al consumo antes de medir & Comprende que hay diferencia de rendimiento pero no puede explicar la causa & No distingue las complejidades ni su implicación energética \\
\hline
\textbf{Uso de CodeCarbon} & Instala y ejecuta CodeCarbon; obtiene al menos un par de mediciones válidas; interpreta la salida del script & Ejecuta correctamente con ayuda docente para un valor de n & Instala pero no consigue ejecutar sin apoyo continuado & No instala o los valores son constantes 0 \\
\hline
\textbf{Análisis de resultados} & Tabla comparativa para los 3 valores de n, gráfico de barras, calcula la reducción porcentual e interpreta la escala de producción & Tabla comparativa y porcentaje de reducción correctos, sin gráfico & Datos para un único valor de n, sin cálculo de reducción & Datos incompletos o sin comparación entre implementaciones \\
\hline
\textbf{Criterio Green IT} & Mide con CodeCarbon las 3 mejoras propuestas, documenta la reducción de cada una y extrapola el ahorro a 10 M peticiones/día & Mide 2 de las 3 mejoras con reducción documentada y cálculo de extrapolación & Mide 1 mejora con CodeCarbon; reducción y extrapolación incompletas & No mide con CodeCarbon o los datos son metodológicamente inválidos \\
\hline
\textbf{Informe técnico} & Informe estructurado con hipótesis previa, tabla de datos, perfil CPU, gráfico, conclusión técnica y reflexión sobre producción & Informe con tabla, análisis básico y conclusión, sin perfil CPU o gráfico & Informe parcial: faltan secciones o el análisis es superficial & Entrega ausente o con menos del 50 \% de los apartados \\
\hline
\end{tabular}}
\caption{Rúbrica P1: \textit{¿Cuánto pesa tu algoritmo?}}
\label{tab:rubrica-p1-cuanto-pesa-tu-algoritmo}
\end{table}

\FloatBarrier
\newpage 
\subsection{Rúbrica P2: \textit{El servidor que no para de crecer}}

\begin{table}[htbp]
\centering
\scalebox{0.92}{\footnotesize
\begin{tabular}{|p{2.6cm}|p{2.6cm}|p{2.6cm}|p{2.6cm}|p{2.6cm}|}
\hline
\textbf{Criterio} & \textbf{Excelente (4)} & \textbf{Bien (3)} & \textbf{En desarrollo (2)} & \textbf{Insuficiente (1)} \\
\hline
\textbf{Diagnóstico técnico} & Identifica y cuantifica todas las causas, señalando qué líneas generan la retención & Identifica la causa principal con medición de crecimiento de memoria & Identifica al menos una causa sin apoyo de datos & No identifica la causa del crecimiento \\
\hline
\textbf{Solución y justificación} & Justifica la solución, contrasta con alternativas y argumenta el trade-off asumido & Justifica la solución con datos y menciona el patrón arquitectónico aplicado & Reduce la RAM con una solución funcional, sin argumentar la elección & La refactorización no reduce el crecimiento o no está justificada \\
\hline
\textbf{Datos y medición} & Tabla completa para $n\in\{100,500,1000\}$ + estimación CO$_2$eq + extrapolación a 8~h & Tabla para los 3 valores de n con \% de reducción, sin extrapolación & Tabla antes/después para un solo valor de n & No presenta tabla comparativa \\
\hline
\textbf{Criterio Green IT} & Calcula el carbono incorporado ahorrado por 2 años de vida extra del servidor (P.III); resultado justificado & Realiza el cálculo de carbono incorporado con pequeños errores & Menciona el carbono incorporado pero no lo cuantifica & No conecta la optimización con el principio Hardware Eficiente \\
\hline
\textbf{Reflexión y transferencia} & Transfiere el aprendizaje a un contexto profesional nuevo con argumentos de diseño & Conecta la decisión con principios POO (encapsulamiento, responsabilidad única) & Reflexión genérica sin conexión con POO & Reflexión ausente o de una línea \\
\hline
\end{tabular}}
\caption{Rúbrica P2: \textit{El servidor que no para de crecer}}
\label{tab:rubrica-p2-el-servidor-que-no-para-de-crecer}
\end{table}

\newpage
\FloatBarrier
\subsection{Rúbrica P3: \textit{Datos que cuestan energía}}

\begin{table}[htbp]
\centering
\scalebox{0.85}{\footnotesize
\begin{tabular}{|p{2.6cm}|p{2.6cm}|p{2.6cm}|p{2.6cm}|p{2.6cm}|}
\hline
\textbf{Criterio} & \textbf{Excelente (4)} & \textbf{Bien (3)} & \textbf{En desarrollo (2)} & \textbf{Insuficiente (1)} \\
\hline
\textbf{Corrección funcional} & Ambas estrategias producen exactamente el mismo resultado (diferencia \textless{} 0,001 en la media); el streaming calcula las 4 estadísticas en un único paso & Mismo resultado para media y co2\_ppm; faltan desviación típica o extremos & La media de temperatura es correcta pero co2\_ppm presenta errores & Las medias difieren entre estrategias o el streaming acumula todos los chunks en memoria \\
\hline
\textbf{Medición de tres métricas} & Tabla con CO$_2$, tiempo y memoria pico para ambas estrategias; gráfico de uso de memoria a lo largo del tiempo & Tabla completa sin gráfico & Solo dos de las tres métricas medidas correctamente & Una sola métrica o medición con errores metodológicos \\
\hline
\textbf{Comprensión de la complejidad espacial} & Explica $O(n)$ vs. $O(k)$ con notación correcta; incluye cálculo del tamaño máximo de memoria en función del dataset y el chunk\_size & Explica la diferencia sin usar notación formal; comprensión correcta & Describe la diferencia intuitivamente sin formalización & No relaciona la estrategia con la complejidad espacial \\
\hline
\textbf{Criterio Green IT} & Reducción de memoria pico \textgreater{}80 \% documentada; calcula el coste energético mensual en un servidor AWS t3.medium usando los tiempos medidos & Reducción \textgreater{}60 \% con medición correcta; sin cálculo de coste en cloud & Diferencia \textless{}40 \% (chunk\_size demasiado grande o dataset demasiado pequeño) & No mide memoria o las mediciones no son comparables \\
\hline
\textbf{Extensión: compresión gzip} & Mide el impacto del CSV comprimido (gzip) en tiempo, memoria y CO$_2$; lo incluye como tercera variante en la tabla & Intenta la lectura comprimida y reporta si el impacto es positivo o negativo & Intenta la compresión pero no consigue ejecutarla correctamente & No aborda la extensión \\
\hline
\end{tabular}}
\caption{Rúbrica P3: \textit{Datos que cuestan energía}}
\label{tab:rubrica-p3-datos-que-cuestan-energia}
\end{table}

\newpage
\FloatBarrier
\subsection{Rúbrica P4: \textit{La cola verde}}

\begin{table}[htbp]
\centering
\scalebox{0.88}{\footnotesize
\begin{tabular}{|p{2.6cm}|p{2.6cm}|p{2.6cm}|p{2.6cm}|p{2.6cm}|}
\hline
\textbf{Criterio} & \textbf{Excelente (4)} & \textbf{Bien (3)} & \textbf{En desarrollo (2)} & \textbf{Insuficiente (1)} \\
\hline

\textbf{Implementación de la estructura de datos} & \texttt{heapq} implementado correctamente con desempate FIFO (contador); \texttt{añadir()} y \texttt{ejecutar\_pendientes()} funcionan sin errores & \texttt{heapq} correcto sin desempate FIFO; el orden entre tareas de igual prioridad no está garantizado & Implementación con lista ordenada manualmente en lugar de \texttt{heapq}; funcional pero ineficiente & Cola desordenada o con errores en la gestión de prioridades \\
\hline
\textbf{Integración con Electricity Maps / simulación} & Integración correcta con la API real (o simulación configurable); maneja timeout, error HTTP y fallback a simulación & Integración con la API sin manejo de excepciones & Solo modo simulado sin intentar la integración con la API & Umbral fijo hardcodeado; sin integración externa \\
\hline
\textbf{Lógica de diferimiento} & Las tareas diferidas regresan a la cola con su prioridad original y el informe registra el estado de cada decisión & Las tareas diferidas regresan correctamente pero el informe no registra el estado & Las tareas diferidas se pierden (no regresan a la cola) & Lógica de diferimiento ausente o siempre ejecuta todas las tareas \\
\hline
\textbf{Criterio Green IT} & Demuestra en la presentación en vivo que difiere al menos una tarea cuando la intensidad supera su umbral y la ejecuta en el ciclo siguiente; explica la conexión con el Principio II de la GSE & Difiere correctamente en la demo pero no articula la conexión con Carbon-Aware Computing & Ejecuta todas las tareas independientemente de la intensidad & La intensidad de carbono no se consulta en ningún momento \\
\hline
\textbf{Tests unitarios} & Suite de ≥ 5 tests; todos pasan; incluye test de regresión para el orden de prioridad & Suite de 3-4 tests, todos pasan; falta test de orden de prioridad o historial & Suite de 1-2 tests que solo verifican el caso básico & Sin tests o tests que dependen de la intensidad real de la red \\
\hline
\end{tabular}}
\caption{Rúbrica P4: \textit{La cola verde}}
\label{tab:rubrica-p4-la-cola-verde}
\end{table}

\FloatBarrier

\newpage
\section{Indicadores de impacto real}

Las rúbricas cubren el aprendizaje técnico. Para evaluar si la propuesta genera cambio real se añaden tres indicadores complementarios: (a) cuestionario Likert pre/post (P0 y P4) sobre percepción del impacto ambiental de las TIC y disposición a actuar; (b) comparación de la Ficha C de P0 con una Ficha C de cierre para detectar cambios de hábito auto-reportados; (c) reducción medida de CO$_2$ y memoria RAM entre la versión inicial y la optimizada en P1–P3. Dado el tamaño reducido del grupo y la ausencia de grupo de control, estos indicadores tienen valor orientativo, no experimental.

\section{Viabilidad y recursos necesarios}

La propuesta es viable en cualquier aula de FP con equipos estándar. Todas las herramientas son gratuitas y open-source (Python, CodeCarbon, pandas); el hardware mínimo es un PC con 4 GB de RAM. El coste temporal es de 14 h sobre 135 h del módulo (10 \%), sin desplazar ningún contenido obligatorio. El material completo está disponible en el repositorio GitHub del proyecto. El único riesgo es la falta de internet para la API de P4; la alternativa offline está documentada en el apartado 4.8.

\section{Limitaciones}

La propuesta presenta cuatro limitaciones principales. La \textbf{formación docente}: los guiones son autocontenidos, pero la motivación del profesor influye en cómo el alumnado percibe la relevancia de la temática. La \textbf{medición longitudinal}: el cuestionario pre/post captura cambios inmediatos, no su persistencia a medio plazo. La \textbf{variabilidad de hardware}: las mediciones de CodeCarbon dependen del equipo; las prácticas mitigan esto comparando siempre versiones en el mismo equipo y ajustando el tamaño del dataset (el apartado 4.8). La \textbf{especificidad del contexto}: la propuesta se ha diseñado para DAM 1.º (MP0485) y su traslado a otros ciclos o a ESO/Bachillerato requiere adaptar RA y entregables.
